\section*{Bab \ref{ch:g8:cosine} --- Segitiga Siku-siku dan Kosinus}

\begin{solution}{exo:g8:cosine:1}
\omterm{pb:g8:midpoints:1}{Garis berat} dari \omterm{def:g3:shapes:rightangle}{sudut siku-sikunya}: setengah sisi miringnya, yaitu
$6.5$ cm (\cref{thm:g8:cosine:circle}). \omterm{pb:g6:symmetry:1}{Lingkaran luarnya} berpusat
di titik tengah sisi miringnya dengan \omterm{def:g3:shapes:shapes}{jari-jari} $6.5$ cm.
\end{solution}

\begin{solution}{exo:g8:cosine:2}
Sudut di $A$ siku-siku ($[BC]$ adalah \omterm{def:g4:geometry:circle}{diameter}: bagian 2 pada
\cref{thm:g8:cosine:circle}). Ukurannya memenuhi
$AB^2 + AC^2 \approx 64 = BC^2$, sampai ketelitian pengukurannya.
\end{solution}

\begin{solution}{exo:g8:cosine:3}
Sisi miringnya: $[EF]$ (yang di seberang \omterm{def:g3:shapes:rightangle}{sudut siku-siku} $D$). Sisi
samping $\widehat E$: $[ED]$. Sisi depannya: $[DF]$. Jadi
$\cos\widehat E = \dfrac{ED}{EF}$.
\end{solution}

\begin{solution}{exo:g8:cosine:4}
\emph{1.} sisi sampingnya $= 10 \cos 35^\circ \approx 8.2$.

\emph{2.} sisi miringnya $= \dfrac{6}{\cos 50^\circ} \approx
\dfrac{6}{0.643} \approx 9.3$.
\end{solution}

\begin{solution}{exo:g8:cosine:5}
$\cos\theta = \frac{8}{10} = 0.8$, jadi
$\theta = \cos^{-1}(0.8) \approx 37^\circ$.
\end{solution}

\begin{solution}{exo:g8:cosine:6}
$\cos\theta$ adalah sisi siku-siku dibagi sisi miring, dan sisi
miringnya adalah sisi \emph{terpanjang} pada \omterm{def:g6:shapes:triangles}{segitiga siku-siku}: jadi
\omterm{def:g3:division:remainder}{hasil baginya} selalu lebih kecil daripada $1$. Nilai $1.2$ akan
berarti sisi siku-siku yang lebih panjang daripada sisi miringnya ---
mustahil.
\end{solution}

\begin{solution}{exo:g8:cosine:7}
Rentang mendatarnya $= 25 \cos 12^\circ \approx 25 \times 0.978
\approx 24.5$ m.
\end{solution}

\begin{solution}{exo:g8:cosine:8}
Kosinusnya turun ketika sudutnya membesar
(\cref{prop:g8:cosine:bounds}):
\[
\cos 70^\circ < \cos 45^\circ < \cos 20^\circ .
\]
(Dengan kalkulator: $0.342 < 0.707 < 0.940$.)
\end{solution}

\begin{solution}{exo:g8:cosine:9}
Jarak di tanahnya adalah sisi samping $\theta$:
$\cos 68^\circ = \frac{1.2}{L}$, jadi
$L = \frac{1.2}{0.375} = 3.2$ m. Tingginya (dengan Pythagoras):
$\sqrt{3.2^2 - 1.2^2} = \sqrt{10.24 - 1.44} = \sqrt{8.8} \approx 3.0$
m.
\end{solution}

\begin{solution}{exo:g8:cosine:10}
\emph{1.} $A$ berada pada lingkaran berdiameter $[BC]$: jadi \omterm{def:g3:shapes:rightangle}{sudut
siku-siku} di $A$ (\cref{thm:g8:cosine:circle}).

\emph{2.} $BC = 9$ cm (dua kali \omterm{def:g3:shapes:shapes}{jari-jarinya}). Dengan Pythagoras:
$AC = \sqrt{9^2 - 5.4^2} = \sqrt{81 - 29.16} = \sqrt{51.84} = 7.2$
cm. Dan $\cos\widehat B = \dfrac{AB}{BC} = \dfrac{5.4}{9} = 0.6$.
\end{solution}

\begin{solution}{exo:g8:cosine:11}
Pada \omterm{def:g6:shapes:triangles}{segitiga sama sisi} bersisi $1$, tinggi dari salah satu \omterm{def:g5:solids:def}{titik
sudutnya} jatuh di \emph{titik tengah} sisi yang di seberangnya (yaitu
\omterm{def:g6:symmetry:axis}{sumbu simetrinya}). Setengah segitiganya siku-siku dengan sisi miring
$1$ (satu sisi penuh), sudut $60^\circ$ di \omterm{def:g5:solids:def}{titik sudut} alasnya, dan
sisi samping $\frac12$ (setengah alasnya). Karena itu
\[
\cos 60^\circ = \frac{1/2}{1} = \frac12 .
\]
\end{solution}

\begin{solution}{pb:g8:cosine:1}
\textbf{1.} Pada segitiga $ABH$, yang siku-siku di $H$, \omterm{def:g1:addition:def}{jumlah}
sudutnya (\cref{thm:g7:triangles:sum}) memberi
$\widehat B + 90^\circ + \widehat{BAH} = 180^\circ$, jadi
$\widehat{BAH} = 90^\circ - \widehat B$. Pada segitiga $ABC$,
yang siku-siku di $A$:
$\widehat B + \widehat C + 90^\circ = 180^\circ$, jadi
$\widehat C = 90^\circ - \widehat B$. Karena itu
$\widehat{BAH} = \widehat{ACB}$: kedua segitiga kecilnya mengulang
sudut $\widehat B$, $\widehat C$, $90^\circ$ milik aslinya.

\textbf{2.} Pada $ABC$ (yang siku-siku di $A$), sisi samping
$\widehat B$ adalah $AB$ dan sisi miringnya $BC$:
$\cos\widehat B = \frac{AB}{BC}$. Pada $ABH$ (yang siku-siku di $H$),
sisi samping $\widehat B$ adalah $BH$ dan sisi miringnya
$AB$: $\cos\widehat B = \frac{BH}{AB}$. Perbandingannya hanya
bergantung pada sudutnya (\cref{def:g8:cosine:def}), jadi
\[
\frac{AB}{BC} = \frac{BH}{AB} ;
\]
dengan mengalikan kedua ruasnya dengan $BC \times AB$
(\cref{thm:g8:equations:balance}):
$AB^2 = BH \times BC$.

\textbf{3.} Segitiga $ABC$ (yang siku-siku di $A$) dan $ACH$
(yang siku-siku di $H$) berbagi sudut $\widehat C$. Sisi samping
dibagi sisi miring pada masing-masingnya:
\[
\cos\widehat C = \frac{AC}{CB} = \frac{CH}{AC}
\qquad\Longrightarrow\qquad
AC^2 = CH \times CB .
\]

\textbf{4.} $AB^2 = BH \times BC$ berbunyi $9 = BH \times 5$, jadi
$BH = 1.8$; dan $16 = CH \times 5$ memberi $CH = 3.2$. Periksa:
$BH + HC = 1.8 + 3.2 = 5 = BC$ --- memang harus begitu, karena $H$
terletak pada sisi miringnya antara $B$ dan $C$.

\textbf{5.} Dengan membagi tiap hubungannya dengan $BC$:
\[
BH = \frac{AB^2}{BC},
\qquad
CH = \frac{AC^2}{BC} .
\]
Jadi $BH$ dan $CH$ \omterm{def:g6:propdata:def}{sebanding} dengan $AB^2$ dan $AC^2$: kaki
garis tingginya membelah sisi miringnya dengan perbandingan kuadrat
kedua sisi siku-sikunya ($1.8 : 3.2 = 9 : 16$ pada pertanyaan 4).

\textbf{6.} Dengan menjumlahkan kedua hubungannya lalu memfaktorkan
$BC$ (sifat distributif):
\[
AB^2 + AC^2 = BH \times BC + CH \times BC
= (BH + CH) \times BC = BC \times BC = BC^2 .
\]
Inilah \cref{thm:g8:pythagoras:direct}, yang diperoleh dari kosinusnya
saja --- bukti yang sungguh berbeda dari teka-teki luas empat segitiga
pada \cref{exo:g8:pythagoras:11}.

\textbf{7.} Pythagoras pada $ABH$ (yang siku-siku di $H$, dengan sisi
miring $[AB]$): $AH^2 = AB^2 - BH^2$. Dengan mengganti $AB^2$ dengan
$BH \times BC$ (pertanyaan 2) lalu memfaktorkan $BH$:
\[
AH^2 = BH \times BC - BH^2
= BH \times (BC - BH)
= BH \times HC .
\]

\textbf{8.} Dengan sisi siku-sikunya sebagai alas dan tinggi, luas
$ABC$ adalah $\frac{AB \times AC}{2}$; dengan sisi miringnya $[BC]$
sebagai alas, tingginya persis $AH$, jadi luasnya juga
$\frac{BC \times AH}{2}$ (\cref{thm:g7:areas:triangle}). Dengan
menyamakan lalu mengalikan dua: $AB \times AC = BC \times AH$.

\textbf{9.} $AH = \frac{AB \times AC}{BC} = \frac{3 \times 4}{5}
= 2.4$. Pemeriksaan pertanyaan 7:
$AH^2 = 2.4^2 = 5.76$ dan
$BH \times HC = 1.8 \times 3.2 = 5.76$. Sama.

\textbf{10.} $AH^2 = BH \times HC = 2 \times 8 = 16$, jadi
$AH = \sqrt{16} = 4$ cm.

\textbf{11.} Titik $A$ terletak pada lingkaran berdiameter $[BC]$
(dan berbeda dari $B$ dan $C$), jadi menurut bagian 2 pada
\cref{thm:g8:cosine:circle} segitiga $ABC$ siku-siku di
$A$. Menurut pelukisannya $(AH) \perp (BC)$ dengan $H$ antara $B$ dan
$C$: \omterm{def:g6:lines:objects}{ruas garis} $[AH]$ persis garis tinggi pada pertanyaan 7,
dan $AH^2 = BH \times HC = p \times q$.

\textbf{12.} $AH^2 = 1 \times 2 = 2$: gambarnya melukis, hanya dengan
penggaris dan jangka, \omterm{def:g6:lines:objects}{ruas garis} yang kuadratnya $2$ --- yaitu
diagonal persegi satuan, $\sqrt2 \approx 1.414$, pada
\cref{ex:g8:pythagoras:diagonal}.

\textbf{13.} Resepnya: tatalah ujung ke ujung dua \omterm{def:g6:lines:objects}{ruas garis}
sepanjang $p = 1$ dan $q = n$; gambarlah setengah lingkaran yang
\omterm{def:g4:geometry:circle}{diameternya} seluruh ruas garisnya (panjangnya $n + 1$); dirikanlah
\omterm{def:g6:lines:perp}{garis tegak lurus} di titik sambungannya; garis itu memotong setengah
lingkarannya di titik yang jaraknya ke titik sambungan adalah
$\sqrt{1 \times n} = \sqrt n$. Untuk $\sqrt5$: ambillah $p = 1$ dan
$q = 5$ (setengah lingkaran berdiameter $6$).

\textbf{14.} $A$ berada pada lingkaran berpusat $O$, jadi
$OA = \frac{p+q}{2}$, yaitu \omterm{def:g3:shapes:shapes}{jari-jarinya}. Kalau $H \neq O$, segitiga
$AHO$ siku-siku di $H$ dengan sisi miring $[OA]$, dan sisi siku-siku
lebih pendek daripada sisi miringnya: $AH < OA$. Kalau $H = O$, maka
$AH = OA$ persis. Pada semua keadaannya $AH \leq \frac{p+q}{2}$,
dengan kesamaannya persis ketika $H$ adalah pusatnya --- yaitu ketika
$p = q$.

\textbf{15.} Dengan mengkuadratkan $AH \leq \frac{p+q}{2}$ (kedua
ruasnya positif) lalu memakai $AH^2 = pq$:
\[
p \times q \leq \left(\frac{p+q}{2}\right)^2 ,
\]
dengan kesamaannya persis untuk $p = q$. Bukti ulang secara aljabar,
dengan kesamaan istimewa (\cref{pb:g8:equations:1}):
\[
\left(\frac{p+q}{2}\right)^2 - pq
= \frac{p^2 + 2pq + q^2 - 4pq}{4}
= \frac{p^2 - 2pq + q^2}{4}
= \left(\frac{p-q}{2}\right)^2 ,
\]
yaitu sebuah kuadrat, jadi tidak pernah negatif --- dan bernilai \omterm{def:g1:counting:zero}{nol}
persis ketika $p = q$. \omterm{pb:g8:cosine:1}{Rata-rata geometri} dua bilangan tidak pernah
melampaui setengah \omterm{def:g1:addition:def}{jumlahnya}.

\end{solution}
